\section{Limitations}

The contracts, security-effect projection, and authorization envelopes remain assumptions that must be validated, not automatically trusted outputs of an LLM. Counterfactual execution can falsify a candidate on tested interventions; absent exhaustive finite-domain analysis, it cannot prove global causal-abstraction soundness. E77 retains changed fields conservatively but conflates arbitrary state/output differences with security effects and misses interaction-only effects. State-dependent, asynchronous, callback, remote, and malicious-tool effects may fall outside a sandbox oracle. Its LLM permission plan can also propose unsupported literals or resolvers unless bounded by an independent manifest, and omitted nonempty defaults require totalization before checking. Canonicalization, provenance, and authorization context may be incomplete or stale. Completed long-horizon and adaptive rows cover only their specified strategies. The conditional theorem excludes compromised tools or hosts and requires check--use integrity not established for concurrent remote services.
